\documentclass[../../../../main.tex]{subfiles}
\begin{document}

平面上に2定点 $A$，$B$ があり，線分 $AB$ の長さ $\overline{AB}$ は $2(\sqrt3+1)$ である．
この平面上を動く3点 $P$，$Q$，$R$ があって，つねに
\begin{align*}
  \overline{AP}=\overline{PQ}=2, \hspace{1\zw} \overline{QR}=\overline{RB}=\sqrt2
\end{align*}
なる長さを保ちながら動いている．
このとき，点 $Q$ が動きうる範囲を図示し，その面積を求めよ．

\begin{figure}[htbp]
  \centering
  \begin{tikzpicture}[scale=1.2]
    \coordinate (A) at (0, 0);
    \coordinate (P) at (1.2, -0.7);
    \coordinate (Q) at (2.4, 0.8);
    \coordinate (R) at (3.8, 1.1);
    \coordinate (B) at (4.8, -0.2);

    \draw[thick] (A) -- (P) -- (Q) -- (R) -- (B);

    \fill (A) circle (1.5pt) node[below left] {$A$};
    \fill (P) circle (1.5pt) node[below] {$P$};
    \fill (Q) circle (1.5pt) node[above left] {$Q$};
    \fill (R) circle (1.5pt) node[above] {$R$};
    \fill (B) circle (1.5pt) node[below right] {$B$};
  \end{tikzpicture}
\end{figure}

\end{document}
